\subsection{环的同调维数}

定义 2. --- 称 A 的同调维数，并记作 dh(A)，其为 \(\overline{\mathbf{Z}}\) 中整数 \(n\) （存在两个 A-模 M 与 N 满足 \(\operatorname{Ext}_{\mathbf{A}}^{n}(\mathbf{M},\mathbf{N})\neq0\) ）所成集合的上确界 .

我们有 \(\mathrm{dh}(0) = -\infty\) , \(\mathrm{dh(A)} \geqslant 0\) （如果 \(A \neq 0\) ）。下面将看到 \(\mathrm{dh(A)} = 1\) （如果 \(A\) 是主理想环且不是域），并且，若 \(K\) 是一个交换域

\[
\mathrm{dh} (\mathrm{K} [ \mathrm{X} _ {1}, \dots , \mathrm{X} _ {n} ]) = n.
\]

命题 4. --- 设 n 为整数且 n ≥ 0。以下条件等价：(i) dh(A) ≤ n，

\begin{enumerate}
\def\labelenumi{(\roman{enumi})}
\setcounter{enumi}{1}
\tightlist
\item
  对每个 A-模 M，有 \(\mathrm{dp}_{\mathbf{A}}(\mathbf{M}) \leqslant n\) ,
\end{enumerate}

(ii') 对每个有限生成的 A-模 M，有 \(\mathrm{dp}_{\mathbf{A}}(\mathbf{M}) \leqslant n\) ,

\begin{enumerate}
\def\labelenumi{(\roman{enumi})}
\setcounter{enumi}{2}
\tightlist
\item
  对每个正合序列
\end{enumerate}

\[
0 \rightarrow \mathrm{K} \rightarrow \mathrm{P} _ {n - 1} \rightarrow \mathrm{P} _ {n - 2} \rightarrow \dots \rightarrow \mathrm{P} _ {0}
\]

其中 \(P_{i}\) 是投射的，K 是投射的，

\begin{enumerate}
\def\labelenumi{(\roman{enumi})}
\setcounter{enumi}{3}
\tightlist
\item
  对每个正合序列
\end{enumerate}

\[
\mathrm{I} ^ {0} \rightarrow \mathrm{I} ^ {1} \rightarrow \dots \rightarrow \mathrm{I} ^ {n - 1} \rightarrow \mathrm{N} \rightarrow 0
\]

其中 \(\mathrm{I}^{i}\) 都是内射的，N 是内射的，

\begin{enumerate}
\def\labelenumi{(\alph{enumi})}
\setcounter{enumi}{21}
\tightlist
\item
  每个 A-模都有长度不超过 n 的内射分解。
\end{enumerate}

条件 (i)、(ii) 和 (iii) 的等价性由命题 1 得出。我们显然有 (ii) \(\Rightarrow\) (ii')。因此，我们只需证明 (ii') \(\Rightarrow\) (iv) \(\Rightarrow\) (v) \(\Rightarrow\) (i)。

(ii') \(\Rightarrow\) (iv)：使用(iv)中的记号，设 K 为 \(\mathrm{I}^{0} \to \mathrm{I}^{1}\) 的核。根据 X，第 128 页，推论 4，对每个 A-模 M，我们有一个同构 \(\operatorname{Ext}_{\mathbf{A}}^{1}(\mathbf{M}, \mathbf{N}) \to \operatorname{Ext}_{\mathbf{A}}^{n+1}(\mathbf{M}, \mathbf{K})\) ；于是由(ii')可知，对每个有限生成的 A-模 M， \(\operatorname{Ext}_{\mathbf{A}}^{1}(\mathbf{M}, \mathbf{N})\) 都为零。根据 X，第 93 页，命题 11，这意味着 N 是内射的，从而得到(iv)。

\begin{enumerate}
\def\labelenumi{(\roman{enumi})}
\setcounter{enumi}{3}
\tightlist
\item
  \(\Rightarrow\) (v)：设 M 为一个 A-模。将(iv)应用于正合序列
\end{enumerate}

\[
0 \rightarrow \mathrm{M} \rightarrow \mathrm{I} ^ {0} (\mathrm{M}) \rightarrow \mathrm{I} ^ {1} (\mathrm{M}) \rightarrow \dots \rightarrow \mathrm{I} ^ {n - 1} (\mathrm{M}) \rightarrow \mathrm{K} ^ {n - 1} (\mathrm{M}) \rightarrow 0
\]

由 X，第 52 页，我们得出结论： \(\mathbf{K}^{n-1}(\mathbf{M})\) 是内射的，由此得 (v)。

\begin{enumerate}
\def\labelenumi{(\alph{enumi})}
\setcounter{enumi}{21}
\tightlist
\item
  \(\Rightarrow\) (i)：这由 X，第 100 页，定理 1 得出。
\end{enumerate}

注记. --- 1) 如果 \(\mathrm{dh}(\mathbf{A}) \leqslant n < +\infty\) ，则对每个 A-模 M 和每个右 A-模 P 有 \(\operatorname{Tor}_{n+1}^{\mathbf{A}}(\mathbf{P}, \mathbf{M}) = 0\) ，因为 \(\mathrm{dp}_{\mathbf{A}}(\mathbf{M}) \leqslant n\) （参见引理 1）。

\begin{enumerate}
\def\labelenumi{\arabic{enumi})}
\setcounter{enumi}{1}
\tightlist
\item
  为使 dh(A) 有限，必要且充分的是对每个非零 A-模 M，dp\_A(M) 均为有限。这确实由前述内容及 X，第 135 页，推论 1 得出。
\end{enumerate}

推论. --- 设 A 为左诺特环，并设 n 为整数 \textgreater{} 0. 以下条件等价：

\begin{enumerate}
\def\labelenumi{(\roman{enumi})}
\item
  \(\mathrm{dh(A)}\leqslant n,\)
\item
  对于每一对有限生成的 A-模 M 和 N，我们有 \(\operatorname{Ext}_{\mathbf{A}}^{n+1}(\mathbf{M}, \mathbf{N}) = 0\) ,
\item
  对每个有限生成的左 A-模 M 和每个有限生成的右 A-模 P，我们有 \(\operatorname{Tor}_{n+1}^{\mathbf{A}}(\mathbf{P},\mathbf{M}) = 0\) .
\end{enumerate}

这由命题 2 和 4 得出。

注记. --- 由(i)与(iii)的等价性，若 A 是右且左诺特的，则 dh(A) = dh(A°)。此等式在一般情况下不成立（X，第 204 页，习题 20）。

命题 5. --- 设 A 是左诺特的且具有有限同调维数，并令 \(\mathcal{C}_{0}\) (相应地， \(\mathcal{C}\) ) 为有限生成投射 A-模（相应地，有限生成 A-模）的同构类集合。则 Grothendieck 群的典范同态 \(\mathrm{K}(\mathcal{C}_{0}) \to \mathrm{K}(\mathcal{C})\) 是双射。

这由 X，第 137 页，推论得出。

